\section{Neue Ämter und Verfassungsorgane}

\subsection{Verfassungskammer} \label{subsec:verfassungskammer}

Die Verfassungskammer wird als Ausgleich für die Proporzabweichungen durch das mehrheitssichernde Verhältniswahlverfahren gebraucht.
Sie ist der Bundestag in der Zusammensetzung nach dem ersten Wahlgang und die einzige Institution, die die Verfassung ändern kann.
\footnote{Man kann diskutieren, ob der Bundesrat weiterhin allen Verfassungsänderungen zustimmen muss. Ich gehe erst einmal davon aus, dass es so bleibt.}
Wenn die Wahl nach nur einem Wahlgang entschieden wurde, ist die Verfassungskammer mit dem Bundestag identisch.
Andernfalls tritt sie in der Zusammensetzung zusammen, die der Bundestag hätte, wenn es nur einen Wahlgang gegeben hätte.
Ihr gehören dann auch Kandidaten an, die nach dem ersten Wahlgang gewählt worden wären, aber nicht im Bundestag sitzen.
Sie kommt nur für ihre Entscheidungen zusammen.

Details siehe \mvref{subsec:verfassungsaenderungen}.

\subsection{Bundespräsidialrat} \label{subsec:bundespraesidialrat}

Es liegt in der Natur der Sache, dass vom Bürger gewählte Verfassungsorgane besonders anfällig gegen Populismus sind.
Auch nicht vom Bürger gewählte Verfassungsorgane können mit der Zeit unterwandert werden, indem sie nach und nach mit Demokratiefeinden besetzt werden.
Solche Prozesse brauchen aber Zeit.
Deswegen können nicht-gewählte Verfassungsorgane eine Zeit lang dem Populismus Grenzen setzen.

Der Gedanke ist, den Bürger die gestaltenden Verfassungsorgane (Exekutive und Legislative) selbst wählen zu lassen,
deren Arbeit aber von nicht-Bürger-gewählten Verfassungsorganen kontrollieren zu lassen.
Die kontrollierenden Verfassungsorgane dürfen selbst nicht gestalten, wie das Land sich entwickelt.
Sie sind ja auch nicht gewählt worden.
Sie sollen aber kontrollieren, dass die gewählten Gestalter ihre Arbeit richtig machen.

Ein kontrollierendes Verfassungsorgan haben wir schon.
Das Bundesverfassungsgericht führt Normenkontrollen durch.
Ihm fehlen aber Durchsetzungsmittel, und es fällt keine staatspolitischen Entscheidungen, sondern stellt nur fest, was recht ist.

Ich schlage vor, das Amt des Bundespräsidenten personell zu verstärken und mit mehr Autorität zu versehen.
Der Bundespräsident sollte einem \enquote{Bundespräsidialrat} \footnote{\enquote{Staatsrat} wäre auch ein passender Begriff. Der Bundespräsident als Staatsoberhaupt steht dem Staatsrat vor.
Der Begriff ist nur schon vorbelastet.} vorstehen,
bestehend aus ihm selbst und sechs anderen \enquote{Bundespräsidialräten}.

Der Bundespräsident leitet die Sitzungen und bestimmt was auf die Tagesordnung kommt, fällt aber ansonsten gleichberechtigt mit den anderen Ratsmitgliedern Entscheidungen.

Die personelle Verstärkung wird gebraucht, um der erhöhten staatspolitischen Verantwortung gerecht zu werden, die diesem Amt zugedacht wird.
Der Bundespräsidialrat soll in mehr Situationen als bisher der Bundespräsident allein den Bundestag auflösen oder den Kanzler absetzen können.
Außerdem sollte er Vetorechte für Exekutiv- und Legislativentscheidungen erhalten.

Der Bundespräsident behält seine repräsentativen Aufgaben.
Die Staatsoberhauptsfunktion soll aber auf sieben Räte inklusive ihm selbst aufgeteilt werden.

Der Bundespräsidialrat ist sozusagen eine \enquote{Metaregierung}.
Er regiert selbst nicht, setzt aber die Regierung im Auftrag des Volkes ein und kann sie bei Bedarf auch wieder absetzen.

\subsubsection{Besetzung des Bundespräsidialrats} \label{subsubsec:besetzung_bundespraesidialrat}

Ein kontrollierendes Verfassungsorgan ist nur so gut wie seine Unabhängigkeit von denen, die es kontrollieren soll.
Wie man parteipolitisch neutrale Organe so besetzt, dass eine Mehrheit sie weder schnell mit eigenen Leuten füllen noch durch Blockade lahmlegen kann, ist ein bekanntes Problem, für das es bewährte Lösungen gibt.\footnote{Zur Theorie solcher Organe z.~B. Bruce Ackerman, \enquote{The New Separation of Powers}, Harvard Law Review 113 (2000), S.~633 ff.; Tarunabh Khaitan, \enquote{Guarantor Institutions}, Asian Journal of Comparative Law 16 (2021), S.~S40 ff.; Mark Tushnet, \enquote{The New Fourth Branch}, Cambridge University Press 2021.
Die Empfehlungen der Venedig-Kommission des Europarats zu qualifizierten Mehrheiten und zur Auflösung von Wahlblockaden sind in CDL-PI(2025)023 zusammengefasst, \url{https://www.coe.int/en/web/venice-commission/-/cdl-pi-2025-023-e}.}
Als Vorbild dient die Besetzung des Bundesverfassungsgerichts nach Art.~93 GG und §§~6, 7 und 7a BVerfGG, die 2024 gerade zum Schutz gegen Populisten teilweise ins Grundgesetz übernommen wurde:
Wahl durch zwei Organe mit Zweidrittelmehrheit, zwölf Jahre Amtszeit ohne Wiederwahl, keine Zugehörigkeit zu Parlamenten oder Regierungen, Weiterführung im Amt bis zur Wahl eines Nachfolgers und Ersatzwahl durch das andere Wahlorgan, wenn die Wahl blockiert wird.
Diese Regeln übernehme ich.
Zusätzlich dürfen die Räte keine Parteiämter bekleiden und nach ihrer Amtszeit für eine Karenzzeit keine politischen Ämter übernehmen.
Im Folgenden beschreibe ich nur, was beim Bundespräsidialrat anders ist.

\runinhead{Wahlorgane}

Der Bundespräsident wird wie bisher von der Bundesversammlung gewählt.
Die sechs übrigen Räte werden je zur Hälfte von der Verfassungskammer und vom Bundesrat gewählt.
Die Verfassungskammer tritt an die Stelle des Bundestags, weil eine Regierungsmehrheit dort durch die Mehrheitssicherung leichter eine Zweidrittelmehrheit erreichen könnte.
Die Räte dürfen auch der Verfassungskammer nicht angehören.

\runinhead{Staffelung}

Alle zwei Jahre wird genau ein Sitz neu besetzt, abwechselnd durch Verfassungskammer und Bundesrat.
Bei der Ersteinrichtung werden die Amtszeiten dafür einmalig auf zwei, vier, sechs, acht, zehn und zwölf Jahre festgelegt.
Eine Mehrheit kann so in einer Legislaturperiode höchstens zwei der sechs Sitze neu besetzen.
Für eine eigene Mehrheit von vier Stimmen bräuchte sie acht Jahre, mit einem eigenen Bundespräsidenten sechs Jahre, und das jeweils mit Zweidrittelmehrheit in Verfassungskammer und Bundesrat.

\runinhead{Tagesordnung}

Das Recht des Bundespräsidenten, die Tagesordnung zu bestimmen, darf nicht dazu führen, dass er unliebsame Themen blockieren kann.
Zwei Räte gemeinsam können deshalb jederzeit eine Sitzung verlangen und einen Punkt auf die Tagesordnung setzen lassen, der dann zeitnah behandelt werden muss.

\runinhead{Ausfall von Räten}

Die Weiterführung im Amt hilft gegen eine Blockade nur, solange die bisherigen Räte ihr Amt noch ausüben können.
Stirbt ein Rat oder tritt er zurück, könnte eine Sperrminorität den Rat nach und nach ausdünnen, bis er handlungsunfähig ist.
Die Ersatzwahl durch das andere Wahlorgan hilft nur, wenn die Blockade auf eines der beiden beschränkt ist.
Außerdem würde jedes vorzeitige Ausscheiden die Staffelung durchbrechen und der gerade amtierenden Mehrheit einen zusätzlichen Sitz verschaffen.

Ich schlage deshalb vor, jeden Rat zusammen mit einem persönlichen Stellvertreter zu wählen, in derselben Wahl und mit derselben Mehrheit.
Fällt der Rat aus, rückt der Stellvertreter für den Rest der Amtszeit nach.
Der Ersatz wurde damit gewählt, als noch Einigkeit bestand, und eine spätere Blockade kann ihn nicht mehr verhindern.

Der Rat bleibt beschlussfähig, solange mindestens fünf Räte im Amt sind.
Eine Entscheidung braucht aber immer vier Stimmen, auch wenn weniger Räte im Amt sind.
Als letzte Notlösung wählen die verbleibenden Räte einen frei gebliebenen Sitz bis zur nächsten erfolgreichen regulären Wahl selbst nach.
Das muss die Ausnahme bleiben, weil sich der Rat sonst selbst ergänzen und von jeder demokratischen Bindung lösen würde.

\runinhead{Abberufung}

Ein Bundespräsidialrat kann nur durch das Bundesverfassungsgericht abberufen werden.
Den Antrag dazu können Verfassungskammer oder Bundesrat jeweils mit Zweidrittelmehrheit stellen, ähnlich wie bei der Präsidentenanklage nach Art.~61 GG.
Gründe für eine Abberufung sind eine vorsätzliche Verletzung des Grundgesetzes oder eines Bundesgesetzes, ein schweres Vergehen im Sinne von Abschnitt~\ref{einzelne_parlamentarier_verstossen_gegen_gesetz}, die dauerhafte Unfähigkeit, das Amt auszuüben, oder wenn der Rat nach seinen Zielen oder seinem Verhalten erkennbar darauf ausgerichtet ist, die Demokratie abzuschaffen oder in eine Ordnung umzubauen, die dem freiheitlichen Grundgedanken des Grundgesetzes zuwiderläuft.

\runinhead{Verankerung im Grundgesetz}

Beim Bundesverfassungsgericht steht die Zweidrittelmehrheit nur im Bundesverfassungsgerichtsgesetz und könnte mit einfacher Mehrheit abgeschafft werden.
Beim Bundespräsidialrat sollten deshalb alle Besetzungsregeln im Grundgesetz stehen, auch die Zweidrittelmehrheit sowie die Stellvertreter- und Ausfallregelung.
Nur Fristen und Verfahrensfragen können einem Bundesgesetz überlassen werden.

\runinhead{Gegenseitige Absicherung und Restrisiko}

Auch wer den Rat beherrscht, kann nicht nach Belieben Bundeskanzler und Bundestag absetzen.
Er braucht dafür entweder den gemessenen Volkswillen (\ref{subsubsubsec:vorgezogene_neuwahlen_auf_volkswunsch_hin}) oder eine stichhaltige Begründung, die das Bundesverfassungsgericht überprüfen kann (\ref{absetzung_durch_bundespraesidialrat}).
Ein Missbrauch setzt also voraus, dass auch das Bundesverfassungsgericht unterwandert ist.
Das unterscheidet den Bundespräsidialrat von nicht gewählten Organen in sogenannten Vormundschaftsdemokratien (\emph{tutelary democracy}) wie Thailand, die gewählte Regierungen nach politischem Gutdünken absetzen.

Schwerer zu kontrollieren ist ein unterwanderter Rat, der einfach nichts tut, also kein Veto einlegt und keine Neuwahlen anordnet.
Gegen eine Mehrheit, die über viele Jahre Zweidrittelmehrheiten in Verfassungskammer und Bundesrat erringt, hilft letztlich keine Besetzungsregel, wie Ungarn und Polen gezeigt haben.
Die Staffelung verschafft dem Volk aber die Zeit, die nach Abschnitt~\ref{sec:haertung_der_demokratie} gebraucht wird, um einen Fehlgriff bei der nächsten Wahl zu korrigieren.

\subsection{Bundesamt für Meinungsforschung}\label{subsec:bundesamt_fuer_meinungsforschung}

Für alle Regierungen ist es wichtig zu wissen, was das Volk denkt.
Demokratische Regierungen bilden da keine Ausnahme.
Für den in \ref{subsec:demokratische_notbremse} gemachten Vorschlag, den Volkswillen bei der Entscheidung für vorgezogene Neuwahlen mitzuberücksichtigen,
ist ein zuverlässiges Stimmungsbild der Bevölkerung essentiell.
Man kann sich bei solchen Entscheidungen nicht allein auf Umfragen privater Meinungsforschungsinstitute verlassen.

Ich schlage deshalb vor, ein Bundesamt für Meinungsforschung einzurichten, das ständig damit beschäftigt ist, mit nicht-geheimdienstlichen Methoden,
ein anonymisiertes Meinungsbild der Bevölkerung zu bestimmten Fragen zu erstellen und aktuell zu halten.
Es sollte sich dabei auf eine größere Datenbasis stützen als die kostengünstigen repräsentativen Umfragen privater Meinungsforschungsinstitute.
Deren Ergebnisse sollten aber auf jeden Fall zur Gegenkontrolle mitherangezogen werden.

\runinhead{Unabhängigkeit}
Das Bundesamt für Meinungsforschung ist fachlich unabhängig, ähnlich wie das Statistische Bundesamt.
Das Amt legt seine Methodik offen.

\runinhead{Auftraggeber}
Umfragen in Auftrag geben können die Bundesregierung, die Fraktionen des Bundestags und der Bundespräsidialrat.
Regierung und Opposition müssen dabei gleichermaßen Zugang zum Amt haben, nicht nur die Regierung allein.
Da auch ein großes Amt nur eine begrenzte Kapazität hat, braucht es dafür Kontingente.
Die regelmäßige Erhebung des Neuwahlwunsches und die Umfragen des Bundespräsidialrats dürfen davon nicht abhängen.
Der Auftraggeber bestimmt das Thema, das Amt den Wortlaut der Fragen.
Es kann Suggestivfragen ablehnen.

\runinhead{Veröffentlichung}
Der Auftraggeber entscheidet, ob und wann die Ergebnisse seiner Umfrage veröffentlicht werden.
Dafür gibt es zwei Gründe:
\begin{enumerate}[label=\alph*)]
    \item Regierung und Opposition wollen oft wissen, wie das Volk über eine Frage denkt, bevor sie sich positionieren oder ihre Entscheidung kommunizieren.
    \item Veröffentlichte Umfragen können einen Verstärkereffekt auslösen, nach dem Motto: \enquote{Wenn X\% der Bevölkerung das so sehen, dann kann es ja nicht falsch sein.}
    Handelt es sich dabei um populistisches Gedankengut, dem man entgegenwirken möchte, macht man sich damit das Leben schwerer.
    Zahlen eines staatlichen Amtes haben dabei mehr Gewicht als die eines privaten Instituts.
\end{enumerate}

Damit sich der Auftraggeber nicht die passenden Zahlen herauspicken kann, veröffentlicht das Amt selbst, und zwar vollständig mit dem Wortlaut der Fragen, allen Antworten und der Methodik.

Immer veröffentlicht werden der regelmäßig erhobene Neuwahlwunsch und alle Ergebnisse, auf die der Bundespräsidialrat seine Entscheidungen stützt.
Eine Geheimhaltung würde hier ohnehin nicht halten, da zu viele Menschen die Zahlen kennen.
Schon das Schweigen wäre ein Signal, und es würde Verschwörungstheorien Nahrung geben.

Spätestens nach Ende der folgenden Legislaturperiode werden alle Ergebnisse veröffentlicht.
Strategisch spielen sie dann keine Rolle mehr, aber man kann nachprüfen, ob das Amt missbraucht wurde.
